\documentclass[a4paper,10pt]{article}
\usepackage[utf8]{inputenc}
\usepackage{graphicx}
\usepackage{subfigure}
\usepackage{fancyhdr}
\usepackage{listings}
\usepackage{ gensymb }
\usepackage{hyperref}
\usepackage{multicol}


%opening

\pagestyle{fancy} %eigener Seitenstil
\fancyhf{} %alle Kopf- und Fußzeilenfelder bereinigen
\fancyhead[L]{SEP Abgabegespräch SS2015/16} %Kopfzeile links
\fancyhead[C]{} %zentrierte Kopfzeile
\fancyhead[R]{\includegraphics[width=1cm]{../../template/tugraz_bw}} %Kopfzeile rechts
\renewcommand{\headrulewidth}{0.4pt} %obere Trennlinie
\fancyfoot[C]{} %Seitennummer
\renewcommand{\footrulewidth}{0.4pt} %untere Trennlinie


\lstset{language=C++, basicstyle=\ttfamily\tiny}  

\begin{document}

\section*{TicTacToe}

Entwickeln Sie eine vereinfachte Version des Spieles TicTacToe\\ (Bsp.: \url{http://playtictactoe.org/}).
Es soll ein zweidimensionales Spielfeld angelegt werden. Das Spielfeld ist dabei
3x3 Felder groß. Die Ausgabe des Spielfelds soll wie in der unten gezeigten Konsolenausgabe erfolgen.\\
Das Spiel soll gegen den Computer gespielt werden können, wobei die KI die Token an zufällige Positionen setzen soll.\\
Der Spieler besitzt den Token \texttt{X} (ASCII 88), der Computer den Token \texttt{O} (ASCII 79).
Der Spieler beginnt das Spiel.
Mit dem Befehl \texttt{quit} soll das Spiel beendet werden.\\
Es gibt keine Gewinnbedingung. Bei Auftreten von spielinternen Fehlern kann eine Fehlermeldung ausgegeben werden,
es soll aber normal weitergespielt werden können.

\begin{multicols}{2}

\subsection*{Beispiel}
\begin{verbatim}
Input auf der Konsole: 
sep> 0 0
sep> 1 0
sep> 2 0
sep> quit
\end{verbatim}

\subsubsection*{Ausgabe}

\begin{verbatim}
Konsolen Ausgabe:
-------------
|   |   |   |
-------------
|   |   |   |
-------------
|   |   |   |
-------------
-------------
| X |   |   |
-------------
|   |   | O |
-------------
|   |   |   |
-------------
-------------
| X | X | O |
-------------
|   |   | O |
-------------
|   |   |   |
-------------
Error: Cannot set Tile on occupied Field!



\end{verbatim}

\end{multicols}


\subsection*{Aufruf}
\texttt{./tictactoe}

\subsection*{Bewertung}
\begin{itemize}
 \item Dokumentation und Programmierstil: 1 Punkt
 \item Einhalten der OOP Konzepte: 1 Punkt
 \item Struktur \& Korrektheit: 1 Punkt
\end{itemize}

\newpage


\end{document}
